\documentclass{article}
\usepackage[UTF8]{ctex}
\usepackage{amsmath,amssymb,amsfonts}
\usepackage{enumitem}
\usepackage{geometry}
\geometry{a4paper, margin=1in}
% === 错题本排版优化（不参与题目标识） ===
\setlength{\parindent}{0pt}
\setlength{\parskip}{0.35em plus 0.10em minus 0.08em}
\setlength{\emergencystretch}{3em}
\setlist{itemsep=0.25em, topsep=0.40em, parsep=0pt}
\allowdisplaybreaks[2]
% === 错题本排版优化结束 ===
\begin{document}
\section*{西南交通大学2024年高等代数真题}
\section*{161．西南交通大学2024年高等代数试卷（回忆版）}

\subsection{题目1}
1．已知矩阵 \(A=\left(\begin{array}{lll}1 & 0 & 1 \\ 0 & 2 & 0 \\ 1 & 0 & 1\end{array}\right)\) ，求 \(A^{n}-2 A^{n-1}=\) \(\_\_\_\_\)．\\
\(\_\_\_\_\)

\subsection{题目2}
2．已知矩阵 \(B=\left(\begin{array}{lll}0 & 0 & 1 \\ 0 & 1 & 0 \\ 1 & 0 & 0\end{array}\right)\) ，且 \(A\) 与 \(B\) 相似，则 \(r(A-E)+r(A-2 E)=\) \(\_\_\_\_\)

\subsection{题目3}
3．设 \(\alpha_{1}=(*, *, *), \alpha_{2}=(*, *, *), \alpha_{3}=(*, *, *)\) 不能被 \(\beta_{1}=(*, *, *), \beta_{2}=(*, *, *), \beta_{3}=(*, *, a)\) 线性表出，求 \(a\) 的值为 \(\_\_\_\_\)．

\subsection{题目4}
4．设 \(A, B\) 为正交矩阵，且 \(\frac{|A|}{|B|}=-1\) ，则 \(|A+B|=\)

\subsection{题目5}
5．设 \(f(x)=x^{2}+x+1, g(x)=x^{2023}-(x+1)^{4045}\) ，求 \(f\) 和 \(g\) 的最大公因式．

\subsection{题目6}
6．求实二次型 \(f\left(x_{1}, x_{2}, \cdots, x_{n}\right)=\sum_{i=1}^{n} x_{i}^{2}+4 \sum_{1 \leq i<j \leq n} x_{i} x_{j}\) 的秩和符号差．

\subsection{题目7}
7．讨论方程组

\[
\left\{\begin{array}{l}
x_{1}+a x_{2}+x_{3}=2 \\
x_{1}+x_{2}+2 b x_{3}=2 \\
x_{1}+x_{2}-b x_{3}=-1
\end{array}\right.
\]

在什么情况下无解，有唯一解，有无穷多解，并在有解时求解。

\subsection{题目8}
8．设 \(A=\left(\begin{array}{ccc}17 & 0 & 25 \\ 0 & 3 & 0 \\ 9 & 0 & -13\end{array}\right)\) ．求 \(A\) 的若尔当标准形 \(J\) ，并求可逆矩阵 \(P\) ，使得 \(P^{-1} A P=J\) ．

\subsection{题目9}
9．设 \(A=\left(\begin{array}{ccc}2 & -1 & 3 \\ a & 1 & b \\ 4 & c & 6\end{array}\right)\) ，存在秩为 2 的矩阵 \(B\) ，使得 \(B A=O\) ，求 \(A^{10}\) ．

\subsection{题目10}
10．已知 \(\alpha_{1}, \alpha_{2}, \alpha_{3}\) 是欧氏空间 \(V\) 的一组基，其度量矩阵为

\[
\left(\begin{array}{ccc}
2 & -1 & 1 \\
-1 & 3 & -1 \\
1 & -1 & 2
\end{array}\right) \text {. (数据可能有误) }
\]

记 \(W=L\left(\alpha_{1}+\alpha_{2}, \alpha_{2}+\alpha_{3}\right)\) ．

\begin{itemize}
  \item[（1）] 求 \(W\) 的一组标准正交基。
  \item[（2）] 求 \(W^{\perp}\) 的维数和一组基。
\end{itemize}

\subsection{题目11}
11．设 \(A\) 是秩为 \(r\) 的 \(n\) 阶方阵，证明：\(A^{2}=A\) 的充要条件是存在 \(r \times n\) 矩阵 \(B\) 和 \(n \times r\) 矩阵 \(C\) ，使得 \(A=B C\) ，且 \(B C=E_{r}\) 。

\subsection{题目12}
12．定义 \(\mathbb{R}_{0}=\{a \in \mathbb{R} \mid a>0\}\) 上的＂加法＂和＂数乘＂为 ⊕ ，⊗ ，满足

\[
a \oplus b=a b, k \otimes a=a^{k}, \forall a, b \in \mathbb{R}_{0}, k \in \mathbb{R} .
\]

证明 \(\mathbb{R}_{0}\) 关于 ⊕ 和 ⊗ 构成线性空间，并求 \(\mathbb{R}_{0}\) 的一组基和维数。
\end{document}
